\documentclass[11pt,a4paper]{article}

\usepackage[utf8]{inputenc}
\usepackage[T1]{fontenc}
\usepackage{lmodern}
\usepackage{amsmath,amssymb,amsfonts}
\usepackage{booktabs}
\usepackage{hyperref}
\usepackage[margin=2.5cm]{geometry}
\usepackage{graphicx}
\usepackage{xcolor}
\usepackage{authblk}

\hypersetup{
  colorlinks=true,
  linkcolor=blue!60!black,
  citecolor=blue!60!black,
  urlcolor=blue!60!black,
  pdftitle={A Point-Like Prediction for Dark Matter at 12.8 GHz from a QCD-Anchored Vacuum Condensate},
  pdfauthor={Oleg Kirichenko}
}

% allow line breaks inside \path{...} at underscores and slashes
\renewcommand{\UrlBreaks}{\do\_\do\/\do\-}

\title{\textbf{A Point-Like Prediction for Dark Matter at 12.8\,GHz\\
from a QCD-Anchored Vacuum Condensate:\\
Theory and Haloscope Test Protocol}}

\author{Oleg Kirichenko}
\affil{Independent Researcher; \texttt{urevich55@gmail.com}\\
\url{https://github.com/infosave2007/vmf}}

\date{\today}

\begin{document}

\maketitle

\begin{abstract}
The Null-Vector Gravity (NVG) framework, anchored to a single QCD input
$M_{\Omega,0} = 859 \pm 8$~MeV fixed by lattice sigma terms, contains a
pseudo-Goldstone mode $\theta$ of the vacuum condensate $\Phi = W e^{i\theta}$.
If the same decay constant $f_a$ that sets the $\theta$-mode mass also generates
Majorana neutrino masses through a QCD-anomaly seesaw, the neutrino sector fixes
$f_a = 1.07\times 10^{11}$~GeV with no further freedom. This single number then
yields a \emph{point} prediction for dark matter: a coherent pseudo-scalar of mass
$m_\theta = 53.2~\mu\text{eV}$, detectable as a microwave line at
$f_\theta = 12.81$~GHz with photon coupling
$g_{\theta\gamma\gamma} = 2.08\times 10^{-14}$~GeV$^{-1}$ (KSVZ benchmark).
The predicted product $m_\theta f_a$ falls on the QCD axion band within 2.6\%,
and the frequency lies in the currently unexplored gap between the published
ORGAN mass windows ($25$--$26$ and $63$--$67~\mu$eV) --- directly on the lower
edge of the SMASH band ($50$--$200~\mu$eV) that ORGAN has declared as its
roadmap target. We present the derivation, a calibration of the haloscope signal
formula against the published ORGAN signal estimate, and two concrete test
protocols: (A)~a resonant-cavity array and (B)~a \emph{fixed-position} dielectric
booster stack, where the point-frequency prediction eliminates the scanning
mechanism. For a coherent dark-matter fraction of 24\% (misalignment only),
protocol~B reaches $\text{SNR}=5$ in ${\sim}\,18$~h at boost factor $\beta=10^4$.
Discovery and exclusion criteria are preregistered. All numerical values are
reproducible from the accompanying scripts.

\medskip
\noindent\textbf{Keywords:} axion dark matter, haloscope, dielectric haloscope,
pseudo-Goldstone boson, vacuum condensate, neutrino mass, point prediction
\end{abstract}

\tableofcontents

%======================================================================
\section{Introduction}
\label{sec:intro}

Direct searches for QCD axion dark matter have so far produced exclusion limits
in disjoint mass windows: ADMX covers ${\sim}\,2.7$--$4.2~\mu$eV
\cite{ADMX2018}, HAYSTAC covers ${\sim}\,17$--$20~\mu$eV
\cite{HAYSTAC2025}, ORGAN has scanned $25.45$--$26.27~\mu$eV
\cite{ORGAN26}, $63.2$--$67.1~\mu$eV \cite{ORGAN2022} and
$107.42$--$111.93~\mu$eV \cite{ORGAN1b}, while MADMAX is under construction
for the $40$--$400~\mu$eV range \cite{MADMAX2019}. The intermediate region
around $50$--$60~\mu$eV ($12$--$15$~GHz) remains unexplored.

This note reports that the NVG/VMF framework \cite{NVGrepo}, anchored to a
single QCD input \cite{VMFanchor}, yields a parameter-free prediction
\emph{exactly} in that gap: a dark-matter pseudo-scalar at
$53.2~\mu$eV. Because the prediction is a \emph{point} in $(m, g_{\gamma})$
space rather than a band, the experimental problem simplifies drastically:
no scanning, no tuning mechanism --- a fixed-frequency integration.

Section~\ref{sec:theory} gives the derivation. Section~\ref{sec:signature}
states the observable signature. Section~\ref{sec:status} documents the gap
in published searches. Section~\ref{sec:protocol} defines the test protocol,
including a calibration of the signal formula against published data and
preregistered discovery/exclusion criteria. Section~\ref{sec:caveats}
lists the framework assumptions honestly.

%======================================================================
\section{Theoretical Framework}
\label{sec:theory}

\subsection{The anchor: $M_{\Omega,0}$ from lattice sigma terms}

The framework is anchored to one dimensionful QCD input. Separating the nucleon
mass into current-quark and nonperturbative components via the Feynman--Hellmann
sigma terms ($\sigma_{\pi N} \approx 44$~MeV, $\sigma_{sN} \approx 30$~MeV),
the vacuum value of the nonperturbative component is fixed at
\cite{VMFanchor}
\begin{equation}
M_{\Omega,0} = 859 \pm 8~\text{MeV},
\end{equation}
which sets the condensate energy density scale
$\rho_c \approx 7.09\times 10^{4}$~MeV/fm$^{3}$ used throughout the
construction. No cosmological or astrophysical input is used at any stage.

\subsection{The $\theta$ pseudo-Goldstone mode}

The NVG vacuum condensate is parametrized as $\Phi = W\,e^{i\theta}$, where the
phase $\theta$ is a pseudo-Goldstone mode with the standard anomaly-generated
potential. Its mass obeys the same relation as the QCD axion,
\begin{equation}
m_\theta = \frac{\sqrt{\chi_{\mathrm{top}}}}{f_a},
\label{eq:mtheta}
\end{equation}
with the topological susceptibility $\chi_{\mathrm{top}}^{1/4} = 75.5$~MeV from
lattice QCD \cite{QCDaxionPrecisely}. In the normalization
$m_a(f_a = 10^{12}~\text{GeV}) = 5.691~\mu$eV \cite{QCDaxionPrecisely}, this
becomes $m_\theta = 5.691~\mu\text{eV}\times (10^{12}~\text{GeV}/f_a)$.

\subsection{The $\theta$-seesaw: fixing $f_a$ from the neutrino sector}
\label{sec:seesaw}

The framework hypothesis that makes the prediction parameter-free is the
\emph{$\theta$-seesaw}: the same mode $\theta$ generates Majorana neutrino
masses without right-handed neutrinos, through the two-loop QCD anomaly diagram
coupling $\theta$ to the lepton current,
\begin{equation}
\mathcal{L}_{\theta\nu} \;=\;
\left(\frac{\alpha_s}{4\pi}\right)^{\!2}\frac{v_{\mathrm{EW}}^{2}}{f_a}\,
\nu_L^{T} C \nu_L + \text{h.c.}
\label{eq:seesaw}
\end{equation}
This gives, for the heaviest neutrino,
\begin{equation}
m_3 = \left(\frac{\alpha_s}{4\pi}\right)^{\!2}\frac{v_{\mathrm{EW}}^{2}}{f_a}.
\label{eq:m3}
\end{equation}
With $\alpha_s(M_Z) = 0.1184$, $v_{\mathrm{EW}} = 246.22$~GeV and the
oscillation-and-cosmology anchor $m_3 = 50.3$~meV (normal ordering; the implied
sum $\sum m_\nu = 59$~meV passes the DESI DR2 bound $\sum m_\nu < 64$~meV
unconditionally), Eq.~\eqref{eq:m3} fixes
\begin{equation}
\boxed{\;f_a = \left(\frac{\alpha_s}{4\pi}\right)^{\!2}
\frac{v_{\mathrm{EW}}^{2}}{m_3} = 1.07\times 10^{11}~\text{GeV}\;}
\label{eq:fa}
\end{equation}
with no remaining freedom. We stress that Eq.~\eqref{eq:seesaw} is a
\emph{framework assumption} (the coefficient $(\alpha_s/4\pi)^2 \approx
8.9\times 10^{-5}$ arises from the two-loop anomaly diagram); it is the single
non-standard ingredient, and it is directly testable (Section~\ref{sec:caveats}).

\subsection{Point prediction for the dark-matter mass}

Inserting \eqref{eq:fa} into \eqref{eq:mtheta}:
\begin{equation}
\boxed{\;m_\theta = 5.691~\mu\text{eV}\times
\frac{10^{12}}{1.07\times 10^{11}} = 53.2~\mu\text{eV}\;}
\end{equation}
or, equivalently, a Compton frequency
\begin{equation}
f_\theta = \frac{m_\theta c^2}{h} = 12.81~\text{GHz}.
\label{eq:freq}
\end{equation}

\paragraph{Consistency check (no tuning).} The two inputs above --- $m_\theta$
from the lattice susceptibility and $f_a$ from the neutrino sector --- are
independent. Their product must land on the QCD axion mass--coupling band
$m_a f_a \simeq m_\pi f_\pi \sqrt{z}/(1+z) \approx 5.82\times 10^{15}$~eV$^2$
\cite{QCDaxionPrecisely}. The framework gives
$m_\theta f_a = 5.67\times 10^{15}$~eV$^2$: \textbf{agreement within 2.6\%},
well inside the lattice uncertainty of the band. This agreement is a
non-trivial consistency test of the construction, not an input.

\subsection{Coupling to photons}

The effective interaction is the standard axion-electrodynamics term
\cite{Sikivie1983}
\begin{equation}
\mathcal{L} \supset -\frac{g_{\theta\gamma\gamma}}{4}\,
\theta\, F_{\mu\nu}\tilde F^{\mu\nu}
= g_{\theta\gamma\gamma}\,\theta\,\mathbf{E}\cdot\mathbf{B},
\qquad
g_{\theta\gamma\gamma} = \frac{\alpha}{2\pi}\,
\frac{|E/N - 1.92|}{f_a}.
\label{eq:gagg}
\end{equation}
For the KSVZ benchmark \cite{KSVZ} ($E/N = 0$) and the DFSZ benchmark
\cite{DFSZ} ($E/N = 8/3$):
\begin{equation}
g_{\theta\gamma\gamma}^{\mathrm{KSVZ}} = 2.08\times 10^{-14}~\text{GeV}^{-1},
\qquad
g_{\theta\gamma\gamma}^{\mathrm{DFSZ}} = 8.10\times 10^{-15}~\text{GeV}^{-1}.
\end{equation}
The model-dependent factor $|E/N - 1.92|$ is the only residual theory
uncertainty on the coupling; the KSVZ value is used as the discovery benchmark.

\subsection{Relic abundance: the honest coherent fraction}
\label{sec:abundance}

For a post-inflationary scenario with the natural initial misalignment
$\theta_i = \pi/\sqrt{3}$, the standard misalignment estimate
\cite{Misalignment} gives
\begin{equation}
\Omega_\theta h^2 \simeq 0.12\,\theta_i^2
\left(\frac{f_a}{10^{12}~\text{GeV}}\right)^{7/6}
= 0.0291,
\end{equation}
i.e.\ \textbf{24\% of the local dark-matter density}
($\Omega_{\mathrm{DM}} h^2 = 0.120$). The remaining density may reside in
topological defects of the condensate sector; it does not contribute to the
coherent haloscope signal. All signal powers below use the conservative coherent
fraction $f_{\mathrm{coh}} = 0.24$; the $f_{\mathrm{coh}}=1$ values are also
quoted where relevant. A larger coherent fraction (e.g.\ from defect decay)
would strengthen the signal proportionally.

\subsection{Fifth force: the pseudoscalar null channel}

The Compton range $\lambda_\theta = \hbar c/m_\theta = 3.72$~mm falls inside
the sub-millimetre window of torsion-balance experiments \cite{AxionReview}.
Because $\theta$ is a \emph{pseudo}scalar, the monopole--monopole Yukawa
channel vanishes identically: unpolarized E\"ot-Wash-class searches must see
\emph{nothing} at this range (null prediction, falsified by any monopole signal),
while the spin-dipole channel with coupling $m_N/f_a$ sits at the edge of
current spin-polarized bounds. This constitutes an independent, low-cost
cross-check of the sector.

%======================================================================
\section{Experimental Signature}
\label{sec:signature}

The complete prediction is summarized in Table~\ref{tab:prediction}.

\begin{table}[h]
\centering
\caption{Parameter-free prediction of the NVG $\theta$-DM sector. No quantity
in the right column is fitted to haloscope data.}
\label{tab:prediction}
\begin{tabular}{lcc}
\toprule
Quantity & Predicted value & Origin \\
\midrule
Decay constant $f_a$ & $1.07\times 10^{11}$~GeV & neutrino sector, Eq.~\eqref{eq:m3} \\
Mass $m_\theta$ & $53.2~\mu$eV & Eq.~\eqref{eq:mtheta}, lattice $\chi_{\mathrm{top}}$ \\
Frequency $f_\theta$ & $12.81$~GHz & $m_\theta c^2/h$ \\
Virial linewidth $Q_a$ & ${\sim}\,10^{6}$ & halo velocity dispersion \\
Line width $\Delta f_a$ & $12.8$~kHz & $f_\theta/Q_a$ \\
$g_{\theta\gamma\gamma}$ (KSVZ) & $2.08\times 10^{-14}$~GeV$^{-1}$ & Eq.~\eqref{eq:gagg} \\
$g_{\theta\gamma\gamma}$ (DFSZ) & $8.10\times 10^{-15}$~GeV$^{-1}$ & Eq.~\eqref{eq:gagg} \\
Coherent DM fraction & $24\%$ & misalignment, Sec.~\ref{sec:abundance} \\
$m_\theta f_a$ vs QCD band & $-2.6\%$ deviation & consistency check \\
\bottomrule
\end{tabular}
\end{table}

The signal is a nearly monochromatic microwave excess at 12.81~GHz with
$Q_a \sim 10^6$ intrinsic quality factor, linearly polarized along the applied
magnetic field, scaling as $B_0^2$, and tracking the sidereal modulation of the
halo wind (a secondary discriminant).

%======================================================================
\section{Status of Published Searches}
\label{sec:status}

Table~\ref{tab:gap} compiles the published haloscope coverage around the
predicted mass.

\begin{table}[h]
\centering
\caption{Published haloscope mass windows around the NVG prediction. The
predicted point ($53.2~\mu$eV) lies in the unexplored gap between the two
ORGAN windows, at the lower edge of the SMASH band ($50$--$200~\mu$eV)
\cite{SMASH} that ORGAN declares as its roadmap target \cite{ORGAN2022}.}
\label{tab:gap}
\begin{tabular}{lccc}
\toprule
Experiment & Mass window ($\mu$eV) & Frequency (GHz) & Best limit on $g_{a\gamma\gamma}$ (GeV$^{-1}$) \\
\midrule
ADMX \cite{ADMX2018} & $2.66$--$3.31$ & $0.64$--$0.80$ & KSVZ reached \\
HAYSTAC II \cite{HAYSTAC2025} & ${\sim}\,17.3$--$19.5$ & $4.18$--$4.71$ & ${\sim}\,10^{-13}$ \\
ORGAN \cite{ORGAN26} & $25.45$--$26.27$ & $6.15$--$6.35$ & $2.8\times 10^{-13}$ \\
\textbf{NVG prediction} & $\mathbf{53.2}$ & $\mathbf{12.81}$ & $\mathbf{2.08\times 10^{-14}}$ \\
ORGAN \cite{ORGAN2022} & $63.2$--$67.1$ & $15.3$--$16.2$ & $3.0\times 10^{-12}$ \\
ORGAN 1b \cite{ORGAN1b} & $107.42$--$111.93$ & $26.0$--$27.1$ & cogenesis ALP excluded \\
MADMAX \cite{MADMAX2019,MADMAXproto} & $40$--$400$ (design) & $10$--$100$ & prototype phase \\
\bottomrule
\end{tabular}
\end{table}

Two remarks. First, no published limit reaches
$g = 2.1\times 10^{-14}$~GeV$^{-1}$ anywhere near $53~\mu$eV: the closest
result (ORGAN at $25.5$--$26.3~\mu$eV) is a factor ${\sim}13$ weaker in
coupling, at a different mass. Second, the ORGAN collaboration's stated goal is
to scan the SMASH band $50$--$200~\mu$eV under a constant magnetic field
\cite{ORGAN2022}; the NVG point sits exactly at the lower edge of that band,
making a targeted single-frequency run a natural addition to their program.

%======================================================================
\section{Test Protocol}
\label{sec:protocol}

\subsection{Signal-power formula and its calibration}
\label{sec:calib}

For a microwave cavity haloscope \cite{Sikivie1983,Egge2023},
\begin{equation}
P_{\mathrm{sig}} = g_{\theta\gamma\gamma}^{2}\,
\frac{\rho_\theta}{m_\theta}\,B_0^{2}\,V\,C\,\min(Q_L, Q_a),
\label{eq:psig}
\end{equation}
where $\rho_\theta = f_{\mathrm{coh}}\,\rho_{\mathrm{DM}}$ with
$\rho_{\mathrm{DM}} = 0.4~\text{GeV/cm}^3$, $V$ is the cavity volume,
$C$ the mode form factor, and $Q_L$ the loaded quality factor.

\textbf{Calibration against published data.} The ORGAN collaboration quotes,
for its Phase~1a apparatus, $P_{\mathrm{sig}} \simeq 2.3\times 10^{-25}$~W
(KSVZ) and $3.1\times 10^{-26}$~W (DFSZ) for a $64~\mu$eV axion
\cite{ORGAN2022}. Two checks:
\begin{enumerate}
\item \emph{Apparatus-independent}: the coupling-coefficient ratio
$(g_{\mathrm{KSVZ}}/g_{\mathrm{DFSZ}})^2 = (0.97/0.36)^2 = 7.26$ versus the
published power ratio $2.3/0.31 = 7.42$ --- agreement at the 2\% level,
confirming the normalization convention of Eq.~\eqref{eq:gagg}.
\item \emph{Absolute}: back-solving Eq.~\eqref{eq:psig} from the published
$2.3\times 10^{-25}$~W implies an apparatus product $V C Q \approx 0.44$~m$^3$
(for $B_0 = 7$~T), consistent with an ORGAN-scale cavity at 15.6~GHz with
$Q \sim 10^{4}$--$10^{5}$. The unit chain of the present calculation is thus
validated against a real haloscope at the factor-of-two level.
\end{enumerate}

For the dielectric booster concept \cite{MADMAX2017,MADMAXtheory,MADMAX2019} the
analogous expression is
\begin{equation}
P_{\mathrm{sig}} = g_{\theta\gamma\gamma}^{2}\,\rho_\theta\,B_0^{2}\,
\frac{A\,\beta}{m_\theta^{2}},
\label{eq:pstack}
\end{equation}
where $A$ is the disk area and $\beta$ the dielectric boost factor; both
Eq.~\eqref{eq:psig} and \eqref{eq:pstack} have been re-derived within the
Lorentz-reciprocity framework \cite{Egge2023,Egge2024exp}.

\subsection{Protocol A: resonant cavity}
\label{sec:cavity}

At $f_\theta = 12.81$~GHz the TM$_{010}$ mode requires a cylinder of radius
$R = x_{01} c/(2\pi f_\theta) = 8.95$~mm. With length $L = 50$~mm,
$V = 12.6$~cm$^3$, form factor $C = 0.69$, $B_0 = 10$~T and a superconducting
(Nb) cavity with baseline $Q_L = 10^{5}$,
\begin{equation}
P_{\mathrm{sig}} = 6.3\times 10^{-25}~\text{W}
\qquad (f_{\mathrm{coh}} = 0.24).
\end{equation}
A single cavity reaches $\text{SNR} = 5$ only after ${\sim}\,48$~years: it is a
pathfinder, not the discovery configuration. The viable cavity route is a
phase-locked array (Table~\ref{tab:arrays}); $Q_L = 10^6$ at 13~GHz is flagged
optimistic (Nb BCS losses scale as $\omega^2$), so the $Q = 10^5$ rows are the
engineering baseline.

\begin{table}[h]
\centering
\caption{Integration time to $\text{SNR}=5$ for cavity arrays
($B_0 = 10$~T, $T_{\mathrm{sys}} = 1$~K, $f_{\mathrm{coh}} = 0.24$,
per-cavity signal from Eq.~\eqref{eq:psig}).}
\label{tab:arrays}
\begin{tabular}{lcccc}
\toprule
& $N = 64$ & $N = 256$ \\
\midrule
$Q_L = 10^{5}$ (baseline) & 276 days & 69 days \\
$Q_L = 10^{6}$ (optimistic) & 0.3 days & 0.1 days \\
\bottomrule
\end{tabular}
\end{table}

\subsection{Protocol B: fixed-position dielectric booster stack}
\label{sec:stack}

The dielectric haloscope \cite{MADMAX2017} replaces the resonant cavity by
${\sim}80$ dielectric disks (LaAlO$_3$, $\varepsilon \approx 24$) in a strong
magnetic field; constructive interference of the emitted photons provides the
boost $\beta$. In a scanning experiment the disk positions must be moved
mechanically --- the technically hardest subsystem, now being commissioned
\cite{MADMAXmech}. \textbf{The point-frequency prediction eliminates this
subsystem entirely}: at a known frequency the inter-disk spacing is fixed, and
the stack becomes a static optical element.

With $D = 1.2$~m disks ($A = 1.13$~m$^2$, the low-mass-end MADMAX aperture),
$B_0 = 10$~T and $f_{\mathrm{coh}} = 0.24$, Eq.~\eqref{eq:pstack} gives
(Table~\ref{tab:stack}). The boost values $\beta = 10^4$--$10^5$ are the design
goals of the MADMAX program \cite{MADMAX2019}; the prototype results are at an
earlier stage \cite{MADMAXproto}, which we state as the principal technical
risk of this option.

\begin{table}[h]
\centering
\caption{Fixed dielectric stack, $B_0 = 10$~T, $T_{\mathrm{sys}} = 1$~K,
readout bandwidth equal to the axion linewidth $\Delta f_a = 12.8$~kHz
(the stack boost band $\gg \Delta f_a$).}
\label{tab:stack}
\begin{tabular}{lccc}
\toprule
Boost $\beta$ & $P_{\mathrm{sig}}$ ($f_{\mathrm{coh}}=0.24$) & $P_{\mathrm{sig}}$ ($f_{\mathrm{coh}}=1$) & $t(\text{SNR}=5)$ \\
\midrule
$10^{4}$ & $3.1\times 10^{-23}$~W & $1.3\times 10^{-22}$~W & 18 h \\
$10^{5}$ & $3.1\times 10^{-22}$~W & $1.3\times 10^{-21}$~W & 11 min \\
\bottomrule
\end{tabular}
\end{table}

Protocol~B is strictly simpler than the cavity array: no cryogenic tuning
mechanism, no phase-locked multi-cavity readout, a single wide-bore magnet,
horn antenna and a near-quantum-limited amplifier --- while the required boost
$\beta$ is already the stated MADMAX design target in a frequency band that
includes 12.81~GHz.

\subsection{Readout, statistics, and preregistered criteria}
\label{sec:stats}

\paragraph{Noise.} $T_{\mathrm{sys}} \approx 1$~K is adopted: millikelvin
physical temperature plus a flux-driven Josephson parametric amplifier. This is
not an assumption but demonstrated technology --- ORGAN's most recent phase
operates exactly this way \cite{ORGAN26}.

\paragraph{Statistical test.} With the Dicke radiometer relation
\begin{equation}
\mathrm{SNR} = \frac{P_{\mathrm{sig}}}{k_B T_{\mathrm{sys}}}
\sqrt{\frac{t}{\Delta f}},
\end{equation}
the \emph{preregistered} search window is $f_\theta \pm 0.5\%$ (covering any
plausible framework uncertainty in $m_3$ and $\chi_{\mathrm{top}}$), resolved
in bins of width $\Delta f_a$. Because the prediction is a point rather than a
band, the look-elsewhere penalty is negligible; no trials factor beyond the
preregistered window is applied.

\paragraph{Discovery criterion.} A $5\sigma$ excess inside the preregistered
window, persisting under independent re-taking, with the two corroborating
checks: (i) power scaling $\propto B_0^2$ when the magnet current is varied;
(ii) consistency of the line shape with the Maxwell--Boltzmann halo profile.
Reproduction at a second apparatus with a different magnet is required before
any claim.

\paragraph{Exclusion criterion.} A null result at KSVZ-level sensitivity,
$g_{\theta\gamma\gamma} \le 2.1\times 10^{-14}$~GeV$^{-1}$ across the
preregistered window, excludes the $\theta$-seesaw dark-matter hypothesis of
Sec.~\ref{sec:seesaw} at $f_{\mathrm{coh}} = 0.24$. Because exclusion curves
assume the full local density, the effective bound on $g$ for a coherent
fraction $f_{\mathrm{coh}}$ is weaker by $1/\sqrt{f_{\mathrm{coh}}}
\approx 2.0$; this rescaling must be applied when interpreting any null.
A null at DFSZ level ($8.1\times 10^{-15}$~GeV$^{-1}$) would be the definitive
test.

\paragraph{Side-channel cross-checks.} The same $f_a$ predicts
$\sum m_\nu = 59$~meV (CMB+BAO, ongoing) and a $0\nu\beta\beta$ window
$m_{\beta\beta} \in [0,\, 6.4]$~meV (LEGEND-1000, nEXO): these provide
framework-level corroboration independent of haloscopes.

%======================================================================
\section{Honest Caveats}
\label{sec:caveats}

\begin{enumerate}
\item \textbf{Framework assumption.} The $\theta$-seesaw, Eq.~\eqref{eq:seesaw},
is a postulate of the NVG construction (motivated by the two-loop anomaly
diagram), not a standard-model result. If Eq.~\eqref{eq:m3} fails --- e.g.\
$m_3$ from a future 0$\nu\beta\beta$/oscillation combination disagrees with
$50.3$~meV --- the entire frequency prediction moves. The haloscope test is
precisely the discriminator.
\item \textbf{Model dependence of the coupling.} The factor
$|E/N - 1.92|$ in Eq.~\eqref{eq:gagg} is benchmarked to KSVZ/DFSZ; a
different UV completion rescales $g_{\theta\gamma\gamma}$ without moving the
frequency. The mass prediction is unaffected.
\item \textbf{Coherent fraction.} All powers use $f_{\mathrm{coh}} = 0.24$
(misalignment only). If most of the sector's density sits in incoherent defects,
haloscope signals weaken accordingly; if defects decay into coherent modes, they
strengthen. This is the dominant astrophysical uncertainty and is why the
exclusion criterion carries the $1/\sqrt{f_{\mathrm{coh}}}$ rescaling.
\item \textbf{Technology risk.} $Q_L = 10^6$ at 13~GHz and $\beta = 10^5$ are
program goals, not demonstrated hardware. The baseline rows ($Q_L = 10^5$,
$\beta = 10^4$) reflect the validated state of the art.
\end{enumerate}

%======================================================================
\section{Conclusion}
\label{sec:concl}

A single neutrino-sector input of the NVG/VMF framework fixes
$f_a = 1.07\times 10^{11}$~GeV, which through lattice-QCD physics yields a
point dark-matter prediction: $m_\theta = 53.2~\mu$eV, a microwave line at
$12.81$~GHz with KSVZ-level photon coupling. The point lands on the QCD axion
band within 2.6\% (consistency, not tuning), sits in the currently unexplored
gap between published ORGAN windows, and is directly on that collaboration's
declared SMASH roadmap. Because the prediction is a point, the search needs no
scanning: a fixed dielectric booster stack reaches discovery-level significance
in under a day of integration at design boost. The protocol above is
preregistered: discovery requires $5\sigma$ plus $B^2$-scaling and
cross-reproduction; exclusion at KSVZ sensitivity across the window closes the
$\theta$-seesaw dark-matter hypothesis. Either outcome is informative --- which
is the property a prediction must have to deserve an experiment.

%======================================================================
\section*{Appendix A: Reproducibility}
\label{app:repro}

All numbers in this note are produced by scripts in the accompanying
repository (\url{https://github.com/infosave2007/vmf},
Zenodo-archived releases):
\begin{sloppypar}
\begin{itemize}
\item \path{generator_project/calc_theta_haloscope_design.py} ---
geometry, signal powers (Eqs.~\eqref{eq:psig}--\eqref{eq:pstack}), ORGAN
calibration, integration times, published-search status;
\item \path{generator_project/calc_theta_haloscope_v2.py} ---
frequency derivation, axion-band consistency check (2.6\%);
\item \path{verification/nvg_neutrino_seesaw.py} --- Eq.~\eqref{eq:m3},
fixing $f_a$; \path{verification/nvg_theta_fifth_force.py} ---
pseudoscalar null channel.
\end{itemize}
\end{sloppypar}

\begin{thebibliography}{99}

\bibitem{VMFanchor}
O.~Kirichenko,
\emph{Lattice Sigma Terms as an Anchor for the Dense Nuclear Matter Equation of State},
Zenodo preprint, doi:10.5281/zenodo.20214457 (2025).

\bibitem{NVGrepo}
O.~Kirichenko,
\emph{NVG-Research: Null-Vector Gravity and Vacuum Mass Fraction framework,
verification suite and scripts},
\url{https://github.com/infosave2007/vmf}.

\bibitem{Weinberg1978}
S.~Weinberg,
\emph{A New Light Boson?},
Phys.\ Rev.\ Lett.\ \textbf{40}, 223 (1978).

\bibitem{Wilczek1978}
F.~Wilczek,
\emph{Problem of Strong $P$ and $T$ Invariance in the Presence of Instantons},
Phys.\ Rev.\ Lett.\ \textbf{40}, 279 (1978).

\bibitem{KSVZ}
J.~E.~Kim,
\emph{Weak-Interaction Singlet and Strong $CP$ Invariance},
Phys.\ Rev.\ Lett.\ \textbf{43}, 103 (1979);
M.~A.~Shifman, A.~I.~Vainshtein, V.~I.~Zakharov,
\emph{Can Confinement Ensure Natural $CP$ Invariance of Strong Interactions?},
Nucl.\ Phys.\ B \textbf{166}, 493 (1980).

\bibitem{DFSZ}
A.~R.~Zhitnitsky,
\emph{On Possible Suppression of the Axion Hadron Interactions},
Sov.\ J.\ Nucl.\ Phys.\ \textbf{31}, 260 (1980);
M.~Dine, W.~Fischler, M.~Srednicki,
\emph{A Simple Solution to the Strong $CP$ Problem with a Harmless Axion},
Phys.\ Lett.\ B \textbf{104}, 199 (1981).

\bibitem{Sikivie1983}
P.~Sikivie,
\emph{Experimental Tests of the Invisible Axion},
Phys.\ Rev.\ Lett.\ \textbf{51}, 1415 (1983).

\bibitem{Misalignment}
J.~Preskill, M.~B.~Wise, F.~Wilczek,
\emph{Cosmology of the Invisible Axion},
Phys.\ Lett.\ B \textbf{120}, 127 (1983);
L.~F.~Abbott, P.~Sikivie,
\emph{A Cosmological Bound on the Invisible Axion},
Phys.\ Lett.\ B \textbf{120}, 133 (1983);
M.~Dine, W.~Fischler,
\emph{The Not-So-Harmless Axion},
Phys.\ Lett.\ B \textbf{120}, 137 (1983).

\bibitem{QCDaxionPrecisely}
G.~Grilli di Cortona, E.~Hardy, J.~Pardo Vega, G.~Villadoro,
\emph{The QCD axion, precisely},
JHEP \textbf{01} (2016) 034, arXiv:1511.02867.

\bibitem{SMASH}
G.~Ballesteros, J.~Redondo, A.~Ringwald, C.~Tamarit,
\emph{Standard Model-Axion-Seesaw-Higgs Portal Inflation. Five problems of
particle physics and cosmology solved in one stroke},
arXiv:1610.01639 (2016).

\bibitem{AxionReview}
I.~G.~Irastorza, J.~Redondo,
\emph{New experimental approaches in the search for axion-like particles},
Prog.\ Part.\ Nucl.\ Phys.\ \textbf{102}, 89 (2018).

\bibitem{AxionExpReview}
P.~W.~Graham, I.~G.~Irastorza, S.~K.~Lamoreaux, A.~Lindner, K.~A.~van Bibber,
\emph{Experimental Searches for the Axion and Axion-Like Particles},
Annu.\ Rev.\ Nucl.\ Part.\ Sci.\ \textbf{65}, 485 (2015).

\bibitem{ADMX2018}
N.~Du et al.\ (ADMX Collaboration),
\emph{A Search for Invisible Axion Dark Matter with the Axion Dark Matter Experiment},
Phys.\ Rev.\ Lett.\ \textbf{120}, 151301 (2018), arXiv:1804.05750.

\bibitem{HAYSTAC2025}
X.~Bai et al.\ (HAYSTAC Collaboration),
\emph{Dark Matter Axion Search with HAYSTAC Phase II},
Phys.\ Rev.\ Lett.\ \textbf{134}, 151006 (2025), arXiv:2409.08998.

\bibitem{ORGAN2022}
A.~Quiskamp et al.\ (ORGAN Collaboration),
\emph{Direct search for dark matter axions excluding ALP cogenesis in the
63- to 67-$\mu$eV range with the ORGAN experiment},
Sci.\ Adv.\ \textbf{8}, eabq3765 (2022).

\bibitem{ORGAN1b}
A.~Quiskamp et al.\ (ORGAN Collaboration),
\emph{Exclusion of Axionlike-Particle Cogenesis Dark Matter in a Mass Window
above 100 $\mu$eV},
Phys.\ Rev.\ Lett.\ \textbf{132}, 031601 (2024), arXiv:2310.00904.

\bibitem{ORGAN26}
A.~P.~Quiskamp et al.\ (ORGAN Collaboration),
\emph{Near-quantum-limited axion dark matter search with the ORGAN experiment
around 26 $\mu$eV},
Phys.\ Rev.\ D \textbf{111}, 095007 (2025), arXiv:2407.18586.

\bibitem{MADMAX2017}
A.~Caldwell et al.\ (MADMAX Working Group),
\emph{Dielectric Haloscopes: A New Way to Detect Axion Dark Matter},
Phys.\ Rev.\ Lett.\ \textbf{118}, 091801 (2017), arXiv:1611.05865.

\bibitem{MADMAXtheory}
A.~J.~Millar, G.~G.~Raffelt, J.~Redondo, F.~D.~Steffen,
\emph{Dielectric Haloscopes to Search for Axion Dark Matter: Theoretical
Foundations},
JCAP \textbf{01} (2017) 061, arXiv:1612.07057.

\bibitem{MADMAX2019}
P.~Brun et al.\ (MADMAX Collaboration),
\emph{A new experimental approach to probe QCD axion dark matter in the mass
range above 40 $\mu$eV},
Eur.\ Phys.\ J.\ C \textbf{79}, 186 (2019), arXiv:1901.07401.

\bibitem{MADMAXsim}
S.~Knirck et al.\ (MADMAX Collaboration),
\emph{Simulating MADMAX in 3D: Requirements for Dielectric Axion Haloscopes},
JCAP \textbf{10} (2021) 034, arXiv:2104.06553.

\bibitem{MADMAXmech}
B.~Ary Dos Santos Garcia et al.\ (MADMAX Collaboration),
\emph{First mechanical realization of a tunable dielectric haloscope for the
MADMAX axion search experiment},
JINST \textbf{19}, T11002 (2024), arXiv:2407.10716.

\bibitem{MADMAXproto}
B.~Ary Dos Santos Garcia et al.\ (MADMAX Collaboration),
\emph{First search for axion dark matter with a MADMAX prototype},
Phys.\ Rev.\ Lett.\ \textbf{135}, 041001 (2025), arXiv:2409.11777.

\bibitem{Egge2023}
J.~Egge,
\emph{Axion haloscope signal power from reciprocity},
JCAP \textbf{04} (2023) 064, arXiv:2211.11503.

\bibitem{Egge2024exp}
J.~Egge et al.,
\emph{Experimental determination of axion signal power of dish antennas and
dielectric haloscopes using the reciprocity approach},
JCAP \textbf{04} (2024) 005, arXiv:2311.13359.

\end{thebibliography}

\end{document}
